% Laget automatisk av elevark/lag_elevark.py fra sim/unnslipningsfart.html. Endringer her blir overskrevet.
\documentclass[11pt]{article}
\usepackage{elevark}
\begin{document}
\elevarkhode{Fysikk 2 · Gravitasjon}{Unnslipningsfart}{https://fysikklab.no/sim/unnslipningsfart.html}{Simuleringen}

Hvor fort må noe skytes opp for at det aldri skal komme ned igjen? Nær bakken kan vi bruke $E_{\mathrm{p}} = mgh$, men langt ute blir tyngdekraften svakere. Da må vi bruke $E_{\mathrm{p}} = -γMm/r$. Skyt rett opp, uten luftmotstand, og se hvor langt prosjektilet kommer.

\section*{Oppgaver}

\begin{oppgaveliste}
\oppgave \fase{gjett}{Gjett}Et prosjektil skytes opp fra jorda med 5~km/s. Gir $E_{\mathrm{p}} = mgh$ for stor eller for liten største høyde?
\svar
\oppgave \fase{test}{Test}Prøv 3, 5, 8 og 10~km/s. Sammenlign største høyde med tilnærmingen mgh. Når er forskjellen liten?
\svar
\oppgave \fase{forklar}{Forklar}Hvorfor er $E_{\mathrm{p}} = -γMm/r$ negativ? Hva betyr det at den samlede energien er negativ, null eller positiv?
\svar
\oppgave \fase{test}{Regn}Vis at unnslipningsfarten er $v = \sqrt{2γM/R}$, og regn den ut for jorda, månen og Mars. Sjekk med simuleringen.
\svar
\oppgave \fase{test}{Regn}Bruk energibevaring og regn ut hvor høyt et prosjektil kommer når det skytes opp fra jorda med 8,0~km/s.
\svar
\oppgave \fase{utfordring}{Utfordring}Månen har ingen atmosfære, mens jorda har det. Bruk unnslipningsfarten og farten til gassmolekyler (noen hundre meter per sekund) til å forklare hvorfor.
\svar
\end{oppgaveliste}

\end{document}
